\wordchapterplain{Appendix A. National Inference and Crop Influence}

The national contrast is linear in log outcomes and can be expressed as \(\widehat{\theta} = \sum_{c,t}^{}a_{ct}\, y_{ct}\). I fit crop and year fixed effects using unexposed observations to obtain residuals \({\widehat{u}}_{ct}\). Treated post-period residuals are then centered by their event-year contrast. The crop-level score and analytic standard error are

\wordeq{0pt}{\psi_{c} = \sum_{t}^{}a_{ct}\,{\widehat{u}}_{ct},\quad\quad SE_{score} = \sqrt{\frac{G}{G - 1}\sum_{c}^{}\psi_{c}^{2}},}{A1}

where \(G\) denotes the number of contributing crop series. For each of 9,999 draws, I draw one multiplier per crop with equal probability from Webb's (2023) six-point distribution:

\wordeq{12pt}{v_{c} \in \left\{ - \sqrt{3/2},\, - 1,\, - \sqrt{1/2},\,\sqrt{1/2},\, 1,\,\sqrt{3/2} \right\}.}{A2}

For each draw, the multiplier statistic is the sum of crop scores multiplied by their respective draws. The interval adds and subtracts the 95th percentile of the statistic's absolute value from the estimated contrast. The p-value is the proportion of draws with an absolute value at least as large as the estimate, with a plus-one adjustment. Area-weighted and event-time statistics use the cell weights implied by equation (6) and the event-time comparisons, retaining the same residuals and draws. The joint pre-pilot test compares the sum of squared standardized contrasts with its distribution across these draws. Independent multipliers preserve dependence within crops but not residual dependence across crops. This score procedure differs from the null-imposed, studentized regression bootstrap used for rice. Webb weights alone do not establish its finite-sample validity.

\begin{wordtable}{A1}{Leave-One-Crop-Out Contrasts}
\begin{wordtabular}[1.25]{L{0.327}C{0.198}C{0.198}C{0.280}}
\wtop
\hdl{0.327}{Omitted crop} & \hd{0.198}{(1)\par Area} & \hd{0.198}{(2)\par Yield} & \hd{0.280}{(3)\par Production} \\
\whead
Soybean & 0.037 & −0.047 & −0.011 \\
Onion & −0.063 & −0.020 & −0.083 \\
Sweet potato & −0.032 & 0.000 & −0.033 \\
Corn & 0.023 & −0.027 & −0.005 \\
Garlic & 0.005 & −0.009 & −0.004 \\
Red pepper & 0.096 & 0.029 & 0.125 \\
\wbottom
\end{wordtabular}

\wordnotes{\emph{Notes}: Log-point contrasts with five treated crops, the full comparison pool, the last pre-pilot reference year, and event years 0-9. Omitted crops are not added to the comparison pool.}
\end{wordtable}
